% !TeX document-id = {0ff8d363-52f6-4b94-8740-bd0bc657a460}
% !TeX program = XeLaTeX
\documentclass[8pt,a4paper]{article}
\usepackage{fullpage}
\usepackage{indentfirst} %首段縮排
\usepackage[top=0.5cm, bottom=0.5cm, left=0.5cm, right=0.5cm]{geometry}
\usepackage{amsmath,amsthm,amsfonts,amssymb,amscd,longtable}
\usepackage{lastpage}
\usepackage{enumerate}
\usepackage{fancyhdr}
\usepackage{mathrsfs}
\usepackage{mathtools}
\usepackage{multicol}
\usepackage{xcolor}
\usepackage{verbatim}
\usepackage{float}  % in order to use H for figure position
\usepackage{graphicx}
\graphicspath{{./Figures_for_main_paper/}{./Figures_for_A_New_Numerical_Method/}}
\DeclareGraphicsExtensions{.pdf,.jpg,.png,.mps} % Portable Document Format,
%\usepackage{luacode}  % This can not let Reference to be showed by Butters
\usepackage{multirow}
\usepackage{makecell}
\usepackage{listings}
\usepackage{multicol}
\usepackage{hyperref}
\usepackage{caption}
\usepackage{xeCJK} %For Traditional Chinese in XeLaTeX
\setCJKmainfont{MingLiU} %For Traditional Chinese in XeLaTeX
\setCJKsansfont{Microsoft JhengHei} %For Traditional Chinese in XeLaTeX

% BibTeX
%\usepackage{bibentry} %Function unknown added by Per Sahlholm
%\usepackage[round,authoryear]{natbib} %Nice author (year) citations
\usepackage[square,numbers]{natbib} %Nice numbered citations
%\usepackage{cite} % Make references as [1-4], not [1,2,3,4]
%\newcommand{\bibspace}{\vspace{2mm}}
\bibliographystyle{unsrtnat} % Unsorted bibliography by Butters
%\bibliographystyle{plainnat} %Choose for including URLs   %Commented by Butters
%\bibliographystyle{unsrt} 
%\bibpunct{[}{]}{,}{y}{,}{,}
%\renewcommand{\bibsection}{} %Removes the References chapter title
%\usepackage[numbib,notlof,notlot,nottoc]{tocbibind} %Adds number to bibliography

\let\ds\displaystyle
\usepackage[shortlabels]{enumitem}
\pagenumbering{gobble}
\hypersetup{
  colorlinks=true,
  linkcolor=blue,
  linkbordercolor={0 0 1}
}
 
\renewcommand\lstlistingname{Algorithm}
\renewcommand\lstlistlistingname{Algorithms}
\def\lstlistingautorefname{Alg.}

\usepackage{listings}
\lstset{language=Matlab,%
    %basicstyle=\color{red},
    breaklines=true,%
    morekeywords={matlab2tikz},
    keywordstyle=\color{blue},%
    morekeywords=[2]{1}, keywordstyle=[2]{\color{black}},
    identifierstyle=\color{black},%
    stringstyle=\color{mylilas},
    commentstyle=\color{mygreen},%
    showstringspaces=false,%without this there will be a symbol in the places where there is a space
    numbers=left,%
    numberstyle={\tiny \color{black}},% size of the numbers
    numbersep=9pt, % this defines how far the numbers are from the text
    emph=[1]{for,end,break},emphstyle=[1]\color{red}, %some words to emphasise
    %emph=[2]{word1,word2}, emphstyle=[2]{style},    
}

\setlength{\parindent}{0.0in}
\setlength{\parskip}{0.05in}

% Edit these as appropriate
\begin{comment}
    \newcommand\course{Digital Control}
    \newcommand\hwnumber{}                  % <-- homework number
    \newcommand\NetIDa{Esteban Roman}           % <-- NetID of person #1
    \newcommand\NetIDb{and his team}           % <-- NetID of person #2 (Comment this line out for problem sets)

    \pagestyle{fancyplain}
    \headheight 35pt
    \lhead{\NetIDa}
    \lhead{\NetIDa\\\NetIDb}                 % <-- Comment this line out for problem sets (make sure you are person #1)
    \chead{\textbf{\Large Final Project \hwnumber}}
    \rhead{\course \\ \today}
    \lfoot{}
    \cfoot{}
    \rfoot{\small\thepage}
    \headsep 1.5em
\end{comment}

\begin{document}
%makecell setting for fonts, spacing etc.

\title{\bf{ My literature review on neural signal decoding articles }}
\author{ 伍錫志 N16074988 }

\begin{comment}
    \setlength\rotheadsize{3cm}
    \renewcommand\theadfont{\small}
    \renewcommand\theadfont{\bfseries}
    \renewcommand\theadset{\renewcommand\arraystretch{0.9}%
    \setlength\extrarowheight{1pt}}
    \renewcommand\cellalign{cc}
    \renewcommand\cellgape{\Gape[2pt]}
\end{comment}

\maketitle       

\section*{Abstract}

Nowadays, digital control has received much attentions for many researchers to investigated before real time digital simulator has been found. 
By utilizing digital signal processing, the signal data can be saved and it is also able to be analyzed for some information. 
Before real-time digital simulator has been established, many researchers tried to develop iteration operator for real time 
approximation of continuous system. It has been verified that the traditional operation prove too slow at that timefor real-time 
simulation due to lack of the numerical integration techniques. Such that, the author in the main paper tried to simplify the iteration discrete 
approximation operator. In this report, the results and comparison has been summarized from two literatures. The results shown that
IBM method could convert the discrete-time system to approximately continuous-time system compare to the other iteration methods.

Digital signal processing has been very developed due to the emerging speed of hardware nowadays.\cite{McLaren1992}
    We can do the simulation approximately real-time as much as an actual event. 
    On the other hand, understanding digital control is not something that able to accomplished. 


\bibliography{./library} \label{sec:references}
\newpage

\end{document}